% =====================================================================
%  Lecture 2026-11-19
%
%  Normalized silently:
%   - B_1(0) and B_r(0) written B(0,1) and B(0,r)
%   - T_lambda = T - lambda I restated for Banach spaces; Section 22
%     introduced it for H
%   - "Ran T_N = N" written dim Ran T_N = N; "ker(W - mu I) < infinity"
%     written dim ker(W - mu I) < infinity
%   - in the diagonal example, the final bound with ||x|| in place of
%     ||x_k||
%
%  CORRECTED SOURCE ERRORS
%   [p1] norm-limit proof: "we show that T x_{n,n} is compact"  ->
%        Cauchy, which is what the estimate shows
%   [p2] weak-to-strong proof: ||y_{n_k} - y|| >= n > 0  ->  >= eta > 0
%   [p3] kernel proposition: ||x_n|| < 1 in the closed unit ball  ->  <=
%   [p3] kernel corollary: the second binomial sum starts at k = 0,
%        producing T^{-1}  ->  k = 1
%
%  ADDITIONS (citations only, on instruction)
%   [p1] Volterra: Theorem 6.3 cited
%   [p1] Lemma 23.4: Lemma 5.3 and Corollary 5.9 cited
%   [p1] Lemma 23.5: Theorem 5.2 cited
%   [p1] Lemma 23.7: the proof of Theorem 5.8 cited for finite-
%        dimensional subspaces being closed, and Theorem 5.4 via
%        Theorem 3.3 for relative compactness
%   [p2] diagonal example: Example 8.8 cited for T bounded, Theorem 4.10
%        for ell_2 Banach
%   [p2] weak to strong: Proposition 19.2 cited for boundedness and
%        uniqueness of weak limits
%   [p2] Corollary 23.12: Lemma 23.4 (2) cited for the contradiction
%   [p3] kernel proposition: Corollary 5.9 cited
%
%  REPORTED BUT UNMODIFIED GAPS
%   [p1] Volterra: only the image of the unit ball is checked, and that
%        Vf is continuous is not stated
%   [p1] Lemma 23.5 (<=): sequences in T(B) are treated, not sequences in
%        its closure
%   [p1] Corollary 23.6: the double subsequence extraction for T_1 + T_2
%        is not written
%   [p1] Theorem 23.8: the bound on ||x_n||, the choice of r, and the
%        convergence of (T_r y_n) are not written
%   [p2] Theorem 23.10: f o T in X' and norm convergence implying weak
%        convergence are not stated
%   [p2] Corollary 23.12: T^{-1} in B(X) uses Lemma 21.7 in its case
%        T in B(X), not proved there
%   [p3] Proposition 23.13: that the limit stays in M is not stated
%   [p3] Theorem 23.15, sigma(V) = {0} and Theorem 23.17 are stated
%        without proof, as in the source; proofs may come with the next
%        lecture
%
%  OMITTED
%   [p1] the two kernel sketches, replaced by one figure on instruction
%   [p1] doodles
% =====================================================================

% [p1]
\section{Compact operators}

\begin{definition}[Compact operator]\label{def:compact-operator}
A linear operator from $(X, \norm{\cdot}_X)$ to $(Y, \norm{\cdot}_Y)$ is
called \emph{compact} if the image of any bounded set is relatively compact
in $Y$.
\end{definition}

\begin{example}[Volterra operator]\label{ex:volterra}
Let $V \colon \Cont{0,1} \to \Cont{0,1}$ with $\norm{\cdot}_\infty$,
$Vf(t) = \int_0^t f(s) \dd s$. The kernel of $V$ is
\[
  K(t,s) =
  \begin{cases}
    1, & s \le t, \\
    0, & s > t,
  \end{cases}
  \qquad
  Vf(t) = \int_0^1 K(t,s)\, f(s) \dd s .
\]
$V$ is compact. Check that the image $V(B)$ of $B = B(0,1)$ is relatively
compact, by the Arzelà--Ascoli theorem, Theorem~\ref{thm:arzela-ascoli}. It is
uniformly bounded: $\norm{f}_\infty \le 1$ gives $\norm{Vf}_\infty \le 1$,
because
$\abs{Vf(t)} = \abs{\int_0^t f(s) \dd s} \le \int_0^t \abs{f(s)} \dd s \le 1$.
It is uniformly equicontinuous: for all $\eps > 0$ there is a $\delta > 0$
such that for all $Vf \in V(B)$ and all $t, t'$ with $\abs{t - t'} < \delta$
we have $\abs{Vf(t) - Vf(t')} < \eps$. Here $\delta = \eps$ works: for
$t' < t$,
\[
  \abs{Vf(t) - Vf(t')} = \Bigl\lvert \int_{t'}^{t} f(s) \dd s \Bigr\rvert
  \le \int_{t'}^{t} \underbrace{\abs{f(s)}}_{\le 1} \dd s
  \le t - t' < \delta = \eps .
\]
\end{example}

\begin{nfigure}{The Volterra operator as a matrix acting on a vector. The kernel
$K(t,s)$ is $1$ on the lower triangle $s \le t$ and $0$ above it, a
continuous analogue of a lower triangular matrix of ones. As in
$(Kf)_t = \sum_s K_{ts} f_s$, the value $Vf(t)$ is row $t$ of $K$ against
$f$, which adds up $f$ over $[0, t]$.}
\begin{tikzpicture}[scale=1.0]
  % the kernel as a matrix: rows t downward, columns s rightward
  \fill[defbg] (0,3) -- (0,0) -- (3,0) -- cycle;
  \draw[rulegrey] (0,0) rectangle (3,3);
  \draw[accent, thick] (0,3) -- (3,0);
  \node[accent] at (0.75,0.85) {$1$};
  \node[rulegrey] at (2.25,2.15) {$0$};
  \node at (1.5,3.35) {$s$};
  \draw[-{Stealth[length=1.6mm]}] (0.2,3.2) -- (2.8,3.2);
  \node at (-0.4,1.5) {$t$};
  \draw[-{Stealth[length=1.6mm]}] (-0.2,2.8) -- (-0.2,0.2);
  % row t
  \draw[thmaccent, line width=2.2pt] (0,1.8) -- (1.2,1.8);
  \draw[thmaccent, dashed] (1.2,1.8) -- (3,1.8);
  \node at (1.5,-0.4) {$K(t,s)$};
  % times
  \node at (3.6,1.5) {$\cdot$};
  % vector f
  \draw[rulegrey] (4.0,0) rectangle (4.5,3);
  \draw[thmaccent, line width=2.2pt] (4.25,3) -- (4.25,1.8);
  \node at (4.25,-0.4) {$f(s)$};
  % equals
  \node at (5.1,1.5) {$=$};
  % vector Vf
  \draw[rulegrey] (5.6,0) rectangle (6.1,3);
  \fill[thmaccent] (5.6,1.72) rectangle (6.1,1.88);
  \node at (5.85,-0.4) {$Vf(t)$};
  \node[thmaccent, right] at (6.15,1.8) {$\displaystyle\int_0^t f(s) \dd s$};
\end{tikzpicture}
\end{nfigure}

\begin{remark}
As a rule of thumb, integral operators are compact. Compact operators have a
simpler spectral theory than bounded linear operators; their spectral theory
is similar to that of matrices.
\end{remark}

% [p1]
\begin{lemma}\label{lem:compact-basic}
\leavevmode
\begin{enumerate}
  \item Every compact operator $T$ is bounded, that is,
    $T \in \Bop{X}{Y}$.
  \item $I \colon X \to X$, $x \mapsto x$, is not compact if
    $\dim X = \infty$, and compact if $\dim X < \infty$.
\end{enumerate}
\end{lemma}

\begin{proof}
(1) Let $B = B(0,1)$. $T(B)$ is relatively compact, so $\cl{T(B)}$ is
compact, then it is bounded, by Lemma~\ref{lem:compact-closed-bounded}, so
$T(B)$ is bounded. Say $\norm{y} \le K$ for all $y \in T(B)$; then
$\norm{T} = \sup_{\norm{x} \le 1} \norm{Tx} \le K$, so $T$ is bounded.

(2) $I(B) = B$, and $B$ is relatively compact if and only if
$\dim X < \infty$, by Corollary~\ref{cor:unit-ball}.
\end{proof}

\begin{lemma}\label{lem:compact-sequential}
$T \in \Bop{X}{Y}$ is compact if and only if for every bounded sequence
$(x_n)$ the image $(Tx_n)$ has a convergent subsequence.
\end{lemma}

\begin{proof}
Assume $T$ is compact and $(x_n)$ is bounded. Then $\cl{\set{Tx_n}}$ is
compact, by definition, so $(Tx_n)$ has a convergent subsequence, by
Theorem~\ref{thm:compact-tfae}.

Conversely, assume that for every bounded $(x_n)$, $(Tx_n)$ has a convergent
subsequence. Let $B \subset X$ be bounded, and consider $T(B)$. $\cl{T(B)}$
is compact, because: take a sequence $(y_n) \subset T(B)$, $y_n = Tx_n$ with
$x_n \in B$. Then $(x_n)$ is bounded, thus $(y_n)$ has a convergent
subsequence, so $\cl{T(B)}$ is compact, by Theorem~\ref{thm:compact-tfae}.
\end{proof}

% [p2]
\begin{corollary}\label{cor:compact-vector-space}
The compact operators form a vector space within $\Bop{X}{Y}$.
\end{corollary}

\begin{proof}
For compact operators $T_1, T_2$, the operators $T_1 + T_2$ and
$\alpha T_1$ are compact, via sequences and Lemma~\ref{lem:compact-sequential}.
\end{proof}

\begin{lemma}\label{lem:finite-rank}
Let $T \in \Bop{X}{Y}$ with $\dim T(X) < \infty$. Then $T$ is compact.
\end{lemma}

\begin{proof}
For every bounded $B \subset X$, $T(B)$ is bounded in $T(X) \subset Y$. As
$\dim T(X) < \infty$, the subspace $T(X)$ is closed in $Y$, as in the proof of
Theorem~\ref{thm:locally-compact}, and its closed bounded sets are compact,
by Theorems~\ref{thm:finite-dim} and~\ref{thm:heine-borel}. So $T(B)$ is
relatively compact.
\end{proof}

\begin{theorem}\label{thm:compact-limit}
Let $(T_n) \subset \Bop{X}{Y}$ be a sequence of compact operators, $Y$ a
Banach space, and $T_n \to T$, that is, $\norm{T - T_n} \to 0$. Then $T$ is
also compact.
\end{theorem}

\begin{proof}
Let $(x_n) \subset X$ be bounded. We need $(Tx_n)$ to have a convergent
subsequence. Use the diagonal method: $(T_1 x_n)$ has a convergent
subsequence $(T_1 x_{1,n})$; $(T_2 x_{1,n})$ has a convergent subsequence
$(T_2 x_{2,n})$; et cetera. Consider the subsequence
$(x_{n,n}) \subset (x_n)$. We show that $(Tx_{n,n})$ is Cauchy. Let
$y_n = x_{n,n}$ and consider
\[
  \norm{Ty_j - Ty_k}
  \le \norm{Ty_j - T_r y_j} + \norm{T_r y_j - T_r y_k}
    + \norm{T_r y_k - T y_k} < \eps .
\]
So $(Ty_n)$ is Cauchy in $Y$, so it converges, as $Y$ is a Banach space.
\end{proof}

\begin{example}\label{ex:diagonal-compact}
Let $T \colon \lp{2} \to \lp{2}$,
$(x_1, x_2, x_3, \dots) \mapsto (\frac{x_1}{1}, \frac{x_2}{2},
\frac{x_3}{3}, \dots)$, bounded by Example~\ref{ex:mult-lp}. $T$ is compact,
because: let
\[
  T_N \colon \lp{2} \to \lp{2}, \qquad (x_1, x_2, \dots) \longmapsto
  \Bigl(\frac{x_1}{1}, \frac{x_2}{2}, \dots, \frac{x_N}{N}, 0, 0, \dots
  \Bigr) .
\]
$T_N$ is compact, since $\dim \operatorname{Ran} T_N = N < \infty$, by
Lemma~\ref{lem:finite-rank}; and $T_N \to T$, so $T$ is compact, by
Theorem~\ref{thm:compact-limit}, $\lp{2}$ being a Banach space by
Theorem~\ref{thm:lp-banach}. Check $T_N \to T$, that is,
$\norm{T - T_N} \to 0$: $\norm{T - T_N} = \sup_{\norm{x} \le 1}
\norm{(T - T_N)x}$, and
\[
  \begin{aligned}
    \norm{(T - T_N)x}
    &= \Bigl\lVert \Bigl(0, \dots, 0, \frac{x_{N+1}}{N+1}, \dots\Bigr)
      \Bigr\rVert_2
    = \Bigl(\sum_{k=N+1}^{\infty} \Bigl\lvert \frac{x_k}{k} \Bigr\rvert^2
      \Bigr)^{1/2} \\
    &\le \frac{1}{N+1} \Bigl(\sum_{k=N+1}^{\infty} \abs{x_k}^2\Bigr)^{1/2}
    \le \frac{1}{N+1} \norm{x} ,
  \end{aligned}
\]
so $\norm{T - T_N} \le \frac{1}{N+1} \to 0$.
\end{example}

\begin{theorem}\label{thm:compact-weak-strong}
If $T \in \Bop{X}{Y}$ is compact, and $(x_n) \subset X$ has
$x_n \xrightarrow{w} x$, then $Tx_n \to Tx$.
\end{theorem}

\begin{proof}
Let $y_n = Tx_n$. As $x_n \xrightarrow{w} x$, $y_n \xrightarrow{w} y = Tx$.
We show that $\norm{y_n - y} \to 0$. Assume $\norm{y_n - y} \not\to 0$, that
is, there is a subsequence $(y_{n_k})$ such that
$\norm{y_{n_k} - y} \ge \eta > 0$. $T$ is compact, and $(x_n)$ is bounded by
Proposition~\ref{prop:weak-limit}, so $(y_{n_k})$ has a convergent
subsequence $\tilde{y}_j \to \tilde{y}$. But then
$\tilde{y}_j \xrightarrow{w} \tilde{y}$, so we must have $y = \tilde{y}$,
since weak limits are unique, by Proposition~\ref{prop:weak-limit}. But this
is a contradiction, since $\norm{\tilde{y}_j - y} \ge \eta$.
\end{proof}

\begin{lemma}\label{lem:compact-ideal}
Let $T \in \Bopone{X}$ be compact, and $S \in \Bopone{X}$. Then $TS$ and $ST$
are also compact.
\end{lemma}

\begin{proof}
$TS$: let $B \subset X$ be bounded; then $S(B)$ is bounded, so $TS(B)$ is
relatively compact.

$ST$: let $(x_n) \subset X$ be bounded. $(Tx_n)$ has a convergent subsequence
$(Tx_{n_k})$; then $(STx_{n_k})$ is also convergent, so $ST$ is compact, by
Lemma~\ref{lem:compact-sequential}, since $(STx_n)$ also has a convergent
subsequence.
\end{proof}

\begin{corollary}\label{cor:compact-zero-spectrum}
Let $X$ be a Banach space. If $\dim X = \infty$, then $0$ is a spectral value
of any compact $T \in \Bopone{X}$.
\end{corollary}

\begin{proof}
Assume $0 \in \rho(T)$; then $T^{-1}$ exists and is bounded, so
$T T^{-1} = I$ must be compact, by Lemma~\ref{lem:compact-ideal}, which is a
contradiction, by Lemma~\ref{lem:compact-basic} (2).
\end{proof}

% [p3]
\begin{proposition}\label{prop:compact-kernel}
Let $X$ be a Banach space, $T \in \Bopone{X}$ compact, and
$T_\lambda = T - \lambda I$. Then for every $\lambda \ne 0$ in $\C$,
$\ker T_\lambda$ is finite-dimensional; that is, if $\lambda$ is an
eigenvalue, its multiplicity is finite.
\end{proposition}

\begin{proof}
Let $M$ be the closed unit ball of $\ker T_\lambda$. We claim that $M$ is
compact, meaning $\dim \ker T_\lambda < \infty$, by
Corollary~\ref{cor:unit-ball}. Let $(x_n) \subset M$; then
$\norm{x_n} \le 1$ and $T_\lambda x_n = 0$, so $x_n = \lambda^{-1} T x_n$.
$T$ is compact, so $(Tx_n)$ has a convergent subsequence $(Tx_{n_k})$; then
$x_{n_k} = \frac{1}{\lambda} T x_{n_k}$ also converges. So $M$ is compact.
\end{proof}

\begin{corollary}\label{cor:compact-kernel-powers}
For $\lambda \ne 0$, $\dim \ker T_\lambda^n < \infty$ for all $n \in \N$,
and
\[
  \set{0} = \ker T_\lambda^0 \subset \ker T_\lambda \subset
  \ker T_\lambda^2 \subset \dots .
\]
\end{corollary}

\begin{proof}
The second part is trivial. For the first,
\[
  T_\lambda^n = (T - \lambda I)^n
  = \sum_{k=0}^{n} \binom{n}{k} T^k (-\lambda)^{n-k}
  = \underbrace{(-\lambda)^n}_{-\mu \ne 0} I
    + \underbrace{T \sum_{k=1}^{n} \binom{n}{k} T^{k-1}
      (-\lambda)^{n-k}}_{W := TS \text{ compact}}
  = W - \mu I ,
\]
where $W$ is compact by Lemma~\ref{lem:compact-ideal}. So
$\dim \ker T_\lambda^n = \dim \ker(W - \mu I) < \infty$, by
Proposition~\ref{prop:compact-kernel}.
\end{proof}

\begin{theorem}\label{thm:compact-eigenvalues}
The set of eigenvalues of a compact operator $T \in \Bopone{X}$ is countable,
and the only possible limit point is $0$; that is, there are finitely many
eigenvalues outside $B(0, r)$ for every $r > 0$.
\end{theorem}

\begin{example}
For the Volterra operator of Example~\ref{ex:volterra},
$Vf(t) = \int_0^t f(s) \dd s$, we have $\sigma(V) = \set{0}$.
\end{example}

\begin{theorem}\label{thm:compact-spectrum}
Let $T \in \Bopone{X}$ be compact on a Banach space $X$. Then every spectral
value of $T$ is an eigenvalue, or zero.
\end{theorem}
